\documentclass[a4,11pt]{aleph-notas}
% Se puede ver la documentación aquí: 
% https://github.com/alephsub0/LaTeX_aleph-notas 

% -- Paquetes adicionales 
\usepackage{enumitem}
\usepackage{array,booktabs,multirow,makecell}
\usepackage{colortbl}
\usepackage{longtable}
\usepackage{url}
\usepackage{pdflscape}
\usepackage{amsmath,amssymb}

% -- Datos 
\institucion{Facultad de Ciencias Exactas, Naturales y Ambientales}
\carrera{Catálogo STEM}
\asignatura{Métodos Numéricos}
\tema{Cronograma y actividades}
\autor{Mario Cueva - Andrés Merino}
\fecha{Periodo 2026-1}

\logouno[0.16\textwidth]{Logos/logoPUCE_04_ac}
\definecolor{colortext}{HTML}{1957d1}
\definecolor{colordef}{HTML}{1957d1}
\fuente{montserrat}


% -- Comandos para tablas
\newcolumntype{C}[1]{>{\hspace{0pt}\centering\arraybackslash}p{#1}}
\newcolumntype{L}[1]{>{\raggedright\arraybackslash}p{#1}}

\definecolor{verde}{RGB}{0, 255, 127}
\definecolor{celeste}{RGB}{68,195,218}

\begin{document}
\addtolength{\headheight}{1.8\baselineskip}
\addtolength{\voffset}{-1.5\baselineskip}

\encabezado

%%%%%%%%%%%%%%%%%%%%%%%%%%%%%%%%%%%%%%%%
\section{Resultados de aprendizaje} 
%%%%%%%%%%%%%%%%%%%%%%%%%%%%%%%%%%%%%%%%

\begin{itemize}[leftmargin=*]
\item 
    \textbf{RdA 1:} Comprender los conceptos fundamentales de los Métodos Numericos, incluyendo la solución de ecuaciones y sistemas, diferenciación e integración numérica y la optimización numerica, destacando su importancia en el análisis y resolución de problemas matemáticos.
    \begin{itemize}[leftmargin=*]
        \item \textbf{Criterio 1.1:}  Describe métodos numéricos para la solución de ecuaciones y sistemas, explicando su funcionamiento y justificando su selección en diferentes contextos matemáticos.
        \item \textbf{Criterio 1.2:}  Comprende los principios de la diferenciación e integración numéricas, destacando su aplicación y relevancia en cálculos matemáticos complejos.
        \item \textbf{Criterio 1.3:} Reconoce las estrategias y técnicas clave de la optimización numérica, analizando su importancia en la resolución eficiente de problemas matemáticos.
    \end{itemize}
\item 
    \textbf{RdA 2:} Computar cálculos relacionadas con los Métodos Numéricos, incluyendo la solución de ecuaciones y sistemas, la diferenciación e integración numéricas, y la optimización numérica, para generar aproximaciones a soluciones.
    \begin{itemize}[leftmargin=*]
        \item \textbf{Criterio 2.1:} Calcula soluciones numéricas de ecuaciones y sistemas con precisión, empleando algoritmos y técnicas apropiadas para diferentes tipos de problemas matemáticos.
        \item \textbf{Criterio 2.2:} Aplica métodos de diferenciación e integración numéricas en una variedad de funciones, demostrando precisión y eficacia en los cálculos realizados.
        \item \textbf{Criterio 2.3:} Emplea técnicas de optimización numérica para solucionar problemas de maximización o minimización, evidenciando comprensión y habilidad en la selección y aplicación de los métodos adecuados.
    \end{itemize}
\item
    \textbf{RdA 3:} Aplicar los conceptos y técnicas de los Métodos Numéricos en la resolución de problemas reales, incluyendo la solución de ecuaciones y sistemas, la diferenciación e integración numéricas, y la optimización numérica, demostrando habilidad para abordar desafíos prácticos y situaciones del mundo real.
    \begin{itemize}[leftmargin=*]
        \item \textbf{Criterio 3.1:} Aplica métodos numéricos en la resolución efectiva de ecuaciones y sistemas complejos, adaptando las técnicas a las características específicas de problemas reales.
        \item \textbf{Criterio 3.2:} Aplica métodos de diferenciación e integración numéricas en la solución de problemas prácticos que requieren análisis detallado y cálculos precisos.
        \item \textbf{Criterio 3.3:} Aplica estrategias de optimización numérica en la resolución de problemas de maximización y minimización en contextos reales, demostrando habilidad en la identificación y aplicación de la técnica más adecuada para cada situación.
    \end{itemize}
\end{itemize}

%%%%%%%%%%%%%%%%%%%%%%%%%%%%%%%%%%%%%%%%
\section{Contenidos generales} 
%%%%%%%%%%%%%%%%%%%%%%%%%%%%%%%%%%%%%%%%

\begin{itemize}
\item 
    Representación de números
\item 
    Definición de errores
\item 
    Resolución de ecuaciones
\item 
    Interpolación polinomial
\item 
    Diferenciación e integración numérica
\item 
    Optimización numérica
\end{itemize}

%%%%%%%%%%%%%%%%%%%%%%%%%%%%%%%%%%%%%%%%
\section{Actividades de evaluación} 
%%%%%%%%%%%%%%%%%%%%%%%%%%%%%%%%%%%%%%%%

% \begin{itemize}[leftmargin=*]
%     \item \textbf{Criterio 1.1}
%         \begin{itemize}[leftmargin=*]
%             \item \textbf{Cuestionario en línea 1 (100\%):} Evaluará la capacidad para identificar cuándo un sistema de ecuaciones lineales es compatible o no, mediante preguntas de opción múltiple y ejercicios gráficos breves.
%         \end{itemize}
        
%     \item \textbf{Criterio 1.2}
%         \begin{itemize}[leftmargin=*]
%             \item \textbf{Cuestionario en línea 2 (100\%):} Medirá la habilidad para distinguir los conceptos básicos relacionados con aplicaciones lineales mediante preguntas conceptuales y análisis de afirmaciones breves.
%         \end{itemize}
        
%     \item \textbf{Criterio 1.3}
%         \begin{itemize}[leftmargin=*]
%             \item \textbf{Cuestionario en línea 3 (100\%):} Evaluará la capacidad para describir los conceptos básicos de la Geometría Analítica en \( \mathbb{R}^2 \), mediante preguntas cortas de desarrollo y ejercicios gráficos breves.
%         \end{itemize}

%     \item \textbf{Criterio 2.1}
%         \begin{itemize}[leftmargin=*]
%             \item \textbf{Examen 1 (100\%):} Evaluará la capacidad para resolver sistemas de ecuaciones lineales utilizando matrices y determinantes mediante ejercicios prácticos de resolución y análisis.
%         \end{itemize}

%     \item \textbf{Criterio 2.2}
%         \begin{itemize}[leftmargin=*]
%             \item \textbf{Examen 2 - parte 1 (100\%):} Medirá la capacidad para resolver operaciones relacionadas con aplicaciones lineales, incluyendo análisis del núcleo e imagen, mediante problemas prácticos y breves justificaciones teóricas.
%         \end{itemize}
    
%     \item \textbf{Criterio 2.3}
%         \begin{itemize}[leftmargin=*]
%             \item \textbf{Examen 2 - parte 2 (100\%):} Evaluará la capacidad para resolver problemas específicos asociados a la Geometría Analítica, especialmente cónicas, mediante ejercicios analíticos y gráficos.
%         \end{itemize}

%     \item \textbf{Criterio 3.1}
%         \begin{itemize}[leftmargin=*]
%             \item \textbf{Reto 1 (100\%):} Consistirá en la modelación de situaciones reales a través de sistemas de ecuaciones lineales, utilizando herramientas computacionales y análisis crítico de resultados.
%         \end{itemize}

%     \item \textbf{Criterio 3.2}
%         \begin{itemize}[leftmargin=*]
%             \item \textbf{Reto 2 - parte 1 (100\%):} Los estudiantes aplicarán transformaciones lineales para resolver problemas concretos de ingeniería, enfatizando el análisis y discusión de resultados obtenidos.
%         \end{itemize}

%     \item \textbf{Criterio 3.3}
%         \begin{itemize}[leftmargin=*]
%             \item \textbf{Reto 2 - parte 2 (100\%):} Evaluará la aplicación práctica de los conceptos de Geometría Analítica en situaciones reales del campo ingenieril, mediante ejercicios integradores donde se usarán técnicas matriciales y vectoriales.
%         \end{itemize}
% \end{itemize}



\begin{landscape}
%%%%%%%%%%%%%%%%%%%%%%%%%%%%%%%%%%%%%%%%
\section{Cronograma de Desarrollo del Curso} 
%%%%%%%%%%%%%%%%%%%%%%%%%%%%%%%%%%%%%%%%

\begin{center}\small
\setlength{\extrarowheight}{0ex}
\setlength{\belowrulesep}{.6ex}
\begin{longtable}{cccL{12cm}L{5cm}}
    \toprule
    &&\thead{Fecha}&\thead{Detalle de contenido} & \thead{Observación} \\
    \midrule
  \endfirsthead
    \multicolumn{5}{l}{\footnotesize \ldots viene de la página anterior}\\
    \toprule
    &&\thead{Fecha}&\thead{Detalle de contenido} & \thead{Observación} \\
    \midrule
  \endhead
        \bottomrule  \multicolumn{5}{r}{\footnotesize Continúa en la siguiente página\ldots}
  \endfoot
        \bottomrule
  \endlastfoot
1	&		&	11-may	&	Introducción a los Métodos Numéricos	&		\\	
	&		&	12-may	&	Representación de números	&		\\	
	&		&	13-may	&	Definición de errores	&		\\	
	&		&	14-may	&	Resolución de ecuaciones: Método de la bisección	&		\\ \midrule	
2	&		&	18-may	&	Resolución de ecuaciones: Método de falsa posición	&		\\	
	&		&	19-may	&	Resolución de ecuaciones: Método de punto fijo	&		\\	
	&		&	20-may	&	Resolución de ecuaciones: Método de Newton–Raphson	&		\\	
	&		&	21-may	&	Resolución de ecuaciones: Método de la secante	&		\\ \midrule	\rowcolor{verde!50}
3	&		&	25-may	&		&	Feriado	\\	
	&		&	26-may	&	Resolución de sistemas de ecuaciones no lineales	&		\\	
	&		&	27-may	&	Resolución de sistemas de ecuaciones lineales	&		\\	
	&		&	28-may	&	Factoriación LU e Inversión de matrices	&		\\ \midrule	
4	&		&	01-jun	&	Métodos iterativos: Jacobi y Gauss–Seidel	&		\\	\rowcolor{celeste!50}
	&		&	02-jun	&	Evaluación escrita	&	Evaluación	\\	
	&		&	03-jun	&	Interpolación polinomial de Newton	&		\\	
	&		&	04-jun	&	Interpolación polinómial de Lagrange	&		\\ \midrule	
5	&		&	08-jun	&	Regresión lineal y polinomial	&		\\	
	&		&	09-jun	&	Diferenciación numérica	&		\\	
	&		&	10-jun	&	Diferenciación numérica de orden superior y parcial	&		\\	
	&		&	11-jun	&	Integración numérica: regla del trapecio	&		\\ \midrule	
6	&		&	15-jun	&	Integración numérica: regla del Simpson	&		\\	
	&		&	16-jun	&	Errores de estimación de diferenciación e integración	&		\\	\rowcolor{celeste!50}
	&		&	17-jun	&	Evaluación escrita	&	Evaluación	\\	
	&		&	18-jun	&	Optimización numérica	&		\\ \midrule	
7	&		&	22-jun	&	Optimización numérica: búsqueda exhaustiva y sección áurea	&		\\	
	&		&	23-jun	&	Optimización numérica: Método de Newton y Brent	&		\\	
	&		&	24-jun	&	Descenso del gradiente I	&		\\	
	&		&	25-jun	&	Descenso del gradiente II	&		\\ \midrule	
8	&		&	29-jun	&	Programación lineal	&		\\	
	&		&	30-jun	&	Optimización numérica con restricciones	&		\\	\rowcolor{celeste!50}
	&		&	01-jul	&	Evaluación escrita	&	Evaluación	\\	
	&		&	02-jul	&	Retroalimentación	&		\\ 
\end{longtable}
\end{center}
\end{landscape}

\end{document} 